\chapter{Введение: от полного процесса к физическому родителю}
\label{ch:v9-introduction}

Том VIII восстановил полный корреляционный оператор, построил
квантово-марковскую динамику, её микроскопические дилатации и условное
межсекторное завершение. Его точная граница состоит не в отсутствии
допустимых форм, а в отсутствии одного происхождения для четырёх независимых
слотов
\begin{equation}
 \mathcal D_{IX}=\{\mathfrak e_{\rm endpoint},E_*,\chi,
 \mathfrak t_{\rm transport}\}.
 \label{eq:v9-intro-four-slots}
\end{equation}
Здесь первая компонента выбирает физическое расширение концевого сектора, вторая задаёт
общий положительный квант энергии, третий --- безразмерную силу связи, а
четвёртый --- фундаментальный закон переноса ячеек.

\section{Центральный вопрос}

Проверяется существование одного функционала, совместимого с вещественной
структурой и калибровочной симметрией,
\begin{equation}
 \mathcal S_{IX}[\Phi;\mathfrak e,E,\chi,\mathfrak t]
 \geqslant \mathcal S_{IX}^{\min},
 \label{eq:v9-intro-parent}
\end{equation}
чьи стационарные условия не просто допускают, а совместно выбирают все
компоненты \eqref{eq:v9-intro-four-slots}. Выбранные значения должны затем
определять полный $44$-канальный процесс без ручной смены носителя и без
подстановки наблюдаемых масс или времён жизни.

\section{Входной контракт}

Программа считается закрытой положительно только при одновременном
выполнении шести требований:
\begin{enumerate}
 \item один общий носитель содержит расширение концевого сектора как внутренний сектор;
 \item одно действие выбирает $E_*>0$ и $\chi$, а не только их допустимый тип;
 \item примитив переноса совместим с GNVW-границей;
 \item едины стационарное состояние и неподвижная алгебра процесса $\mathcal L_{44}$;
 \item полный физический гессиан неотрицателен после факторизации по
 калибровочным направлениям;
 \item получено хотя бы одно слепое безразмерное следствие.
\end{enumerate}
Итоговый счётчик успеха имеет вид
\begin{equation}
 N_{IX}=6/6.
 \label{eq:v9-intro-success-ledger}
\end{equation}

\section{Рабочая и нулевая гипотезы}

Рабочая гипотеза утверждает, что общий носитель и один ограниченный снизу
родительский функционал существуют. Нулевая гипотеза утверждает, что на всех носителях,
допущенных Томом VIII, четыре слота остаются независимыми внешними данными:
\begin{equation}
 N_{\rm inherited\ selection}=0/4.
 \label{eq:v9-intro-inherited-ledger}
\end{equation}
Первый гейт проверяет допустимость этой постановки. Если типы слотов
смешиваются, критерии нельзя вычислить или конструкция требует заранее
подставить целевое наблюдаемое, программа останавливается до построения
нового родительского функционала.

\begin{protocol}[Открытие Тома IX]\begin{enumerate}
\item вход: \textbf{завершённый операторный процесс Тома VIII};
\item неизвестное: \textbf{единый четырёхкомпонентный физический родительский функционал};
\item положительный порог: \textbf{$6/6$};
\item внешняя калибровка: \textbf{запрещена};
\item первый шаг: \textbf{проверка допустимости программы и типов данных}.
\end{enumerate}\end{protocol}